% SPDX-License-Identifier: CC-BY-SA-4.0
% Author: Matthieu Perrin
% Part: <Nom de la partie>
% Section: <Nom de la section>
% Sub-section: <Nom de la sous-section>  % (facultatif, laisser vide si non utilisé)
% Frame: <Titre de la slide>

\begingroup

\newcommand\mcell[5]{
  \draw[draw=#1, fill=#3] #4 +(-.45,-.45) rectangle +(.45,.45);
  \node[#2]            at #4 {#5};
}
\newcommand\mcellN[1]{\mcell{black!50}{}{black!5}{#1}{}}
\newcommand\mcellZ[1]{\mcell{black!50}{black!40}{black!25}{#1}{0}}
\newcommand\mcellA[1]{\mcell{structure!70}{structure}{structure!25}{#1}{1}}
\newcommand\mcellB[1]{\mcell{example!70}{example}{example!25}{#1}{2}}
\newcommand\mcellC[1]{\mcell{alert!70}{alert}{alert!25}{#1}{3}}

\begin{frame}{Jeu du Démineur}

  \on[y=7mm, width=1.1\textwidth]{
    \Probleme{Minesweeper-Consistency}{
      \begin{itemize}
      \item Une grille finie $G \subset \mathbb{Z}^2$
      \item Une fonction partielle $\ell : G \rightharpoonup \{0, \ldots, 8\}$
      \end{itemize}
      
      \begin{center}
        \begin{tikzpicture}[size=3mm]\scriptsize
          
          \mcellZ{(1,9)}
          \mcellZ{(2,9)}
          \mcellZ{(3,9)}
          \mcellZ{(4,9)}
          \mcellZ{(5,9)}
          \mcellZ{(6,9)}
          \mcellZ{(7,9)}
          \mcellZ{(8,9)}
          \mcellZ{(9,9)}

          \mcellZ{(1,8)}
          \mcellZ{(2,8)}
          \mcellZ{(3,8)}
          \mcellZ{(4,8)}
          \mcellZ{(5,8)}
          \mcellZ{(6,8)}
          \mcellA{(7,8)}
          \mcellA{(8,8)}
          \mcellA{(9,8)}

          \mcellZ{(1,7)}
          \mcellZ{(2,7)}
          \mcellZ{(3,7)}
          \mcellZ{(4,7)}
          \mcellA{(5,7)}
          \mcellA{(6,7)}
          \mcellB{(7,7)}
          \mcellN{(8,7)}
          \mcellN{(9,7)}

          \mcellA{(1,6)}
          \mcellA{(2,6)}
          \mcellZ{(3,6)}
          \mcellZ{(4,6)}
          \mcellA{(5,6)}
          \mcellN{(6,6)}
          \mcellB{(7,6)}
          \mcellA{(8,6)}
          \mcellA{(9,6)}

          \mcellN{(1,5)}
          \mcellC{(2,5)}
          \mcellA{(3,5)}
          \mcellA{(4,5)}
          \mcellA{(5,5)}
          \mcellA{(6,5)}
          \mcellA{(7,5)}
          \mcellZ{(8,5)}
          \mcellZ{(9,5)}

          \mcellN{(1,4)}
          \mcellN{(2,4)}
          \mcellN{(3,4)}
          \mcellB{(4,4)}
          \mcellZ{(5,4)}
          \mcellZ{(6,4)}
          \mcellZ{(7,4)}
          \mcellZ{(8,4)}
          \mcellZ{(9,4)}

          \mcellN{(1,3)}
          \mcellN{(2,3)}
          \mcellN{(3,3)}
          \mcellC{(4,3)}
          \mcellZ{(5,3)}
          \mcellZ{(6,3)}
          \mcellZ{(7,3)}
          \mcellA{(8,3)}
          \mcellA{(9,3)}

          \mcellN{(1,2)}
          \mcellN{(2,2)}
          \mcellN{(3,2)}
          \mcellB{(4,2)}
          \mcellA{(5,2)}
          \mcellA{(6,2)}
          \mcellA{(7,2)}
          \mcellA{(8,2)}
          \mcellN{(9,2)}

          \mcellN{(1,1)}
          \mcellN{(2,1)}
          \mcellN{(3,1)}
          \mcellN{(4,1)}
          \mcellN{(5,1)}
          \mcellN{(6,1)}
          \mcellN{(7,1)}
          \mcellN{(8,1)}
          \mcellN{(9,1)}
          
        \end{tikzpicture}
      \end{center}
    }{
      Existe-t-il un ensemble $M \subseteq G$ (positions de mines) tel que
      \begin{itemize}
      \item $M \cap \dom(\ell) = \emptyset$, et
      \item $\forall c \in \dom(\ell), \ell(c) = |\{d \in M \mid d \text{ est adjacent à } c\}|$ ?
      \end{itemize}
    }
  }  

  \onBlock[bottom=3mm]{Théorème (admis)}{
    \textsc{Minesweeper-Consistency} est \textsc{np-}complet.
  }

  \footnoteref{R. Kaye. \emph{Minesweeper is NP-complete}. The Mathematical Intelligencer (2000)}
  
\end{frame}

\endgroup
